\subsection{DEFINITION OF A LIE GROUP}

Let \(G\) be a set. A group structure and an analytic K-manifold structure on \(G\) are called compatible if the following condition holds:

(GL) The mapping \((g, h) \mapsto gh^{-1}\) of \(\mathbf{G} \times \mathbf{G}\) into \(\mathbf{G}\) is analytic.

DEFINITION 1. A Lie group over K is a set G with a group structure and an analytic K-manifold structure such that these two structures are compatible.

A Lie group over \(\mathbf{R}\) (resp. \(\mathbf{C},\mathbf{Q}_p\) ) is called a real (resp. complex, \(p\) -adic) Lie group.

Let \(\mathbf{G}\) be a group with an analytic manifold structure. For \(g, h, g_0, h_0\) in \(\mathbf{G}\) ,

\[
g h ^ {- 1} = (g _ {0} h _ {0} ^ {- 1}) h _ {0} ((g _ {0} ^ {- 1} g) (h _ {0} ^ {- 1} h) ^ {- 1}) h _ {0} ^ {- 1}.\tag{1}
\]

It follows that G is a Lie group if and only if the following three conditions hold:

\((\mathbf{GL}_1)\) for all \(g_0 \in \mathbf{G}\) , the mapping \(g \mapsto g_0 g\) of \(\mathbf{G}\) into \(\mathbf{G}\) is analytic;

\((\mathrm{GL}_2)\) for all \(g_0 \in \mathbf{G}\) , the mapping \(g \mapsto g_0gg_0^{-1}\) of \(\mathbf{G}\) into \(\mathbf{G}\) is analytic in an open neighbourhood of \(e\) ;

\((\mathrm{GL}_3)\) the mapping \((g,h)\mapsto gh^{-1}\) of \(\mathbf{G}\times \mathbf{G}\) into \(\mathbf{G}\) is analytic in an open neighbourhood of \((e,e)\) .

Let G be a Lie group. For all \(g \in G\) , \(\gamma(g)\) and \(\delta(g)\) are automorphisms of the underlying manifold of G. It follows that this manifold is pure (Differentiable and Analytic Manifolds, R, 5.1.7). In particular, the dimension of G at g is equal to dim G for all \(g \in G\) (recall that dim G is an integer \(\geqslant 0\) or \(+\infty\) ).

Since an analytic mapping is continuous, a Lie group is a topological group with the underlying topology of its manifold structure. Let G be a set. A topological group structure and an analytic K-manifold structure on G are called compatible if the group structure and the manifold structure are compatible and the topology on G is the topology underlying the manifold structure.

Lemma 1. Let G be a Lie group, U an open neighbourhood of e, E a complete normed space and \(\phi: \mathrm{U} \to \mathrm{E}\) a chart of the manifold G. There exists a neighbourhood W of e contained in U such that \(\phi \mid W\) is an isomorphism of W (with the right uniform structure) onto \(\phi(W)\) (with the uniform structure induced by that on E).

It can be assumed that \(\phi(e) = 0\) . Let \(U' = \dot{\phi}(U)\) . Let \(\dot{\psi}: U' \to U\) be the inverse mapping of \(\dot{\phi}\) . Let \(V\) be a symmetric open neighbourhood of \(e\) such that \(V^2 \subset U\) and let \(V' = \dot{\phi}(V)\) . We define mappings \(\theta_1, \theta_2\) of \(V' \times V'\) into \(V' \times U'\) as follows:

\[
\begin{array}{l} \theta_ {1} (x, y) = (x, \phi (\psi (x) \psi (y) ^ {- 1})) \\ \theta_ {2} (x, y) = (x, \phi (\psi (y) ^ {- 1} \psi (x))). \end{array}
\]

It is immediately verified that \(\theta_{2}(\theta_{1}(x,y)) = \theta_{1}(\theta_{2}(x,y)) = (x,y)\) for \(x\) , \(y\) sufficiently close to 0. On the other hand, \(\theta_{1}\) and \(\theta_{2}\) are analytic and hence strictly differentiable at \((0,0)\) . Therefore (Differentiable and Analytic Manifolds, R, 1.2.2) there exist a neighbourhood \(W'\) of 0 in \(V'\) and constants \(a > 0\) , \(b > 0\) such that

\[
\begin{array}{r l} a (\| x _ {1} - x _ {2} \| + \| \phi (\psi (x _ {1}) \psi (y _ {1}) ^ {- 1}) - \phi (\psi (x _ {2}) \psi (y _ {2}) ^ {- 1}) \|) \\ & \leqslant \| x _ {1} - x _ {2} \| + \| y _ {1} - y _ {2} \| \\ & \leqslant b (\| x _ {1} - x _ {2} \| + \| \phi (\psi (x _ {1}) \psi (y _ {1}) ^ {- 1}) - \phi (\psi (x _ {2}) \psi (y _ {2}) ^ {- 1}) \|) \end{array}
\]

for all \(x_1, x_2, y_1, y_2\) in \(W'\) . Writing \(x_1 = x_2 = y_2\) , we obtain

\[
a \| \phi (\psi (x _ {1}) \psi (y _ {1}) ^ {- 1}) \| \leqslant \| x _ {1} - y _ {1} \| \leqslant b \| \phi (\psi (x _ {1}) \psi (y _ {1}) ^ {- 1}) \|.\tag{2}
\]

For \(\delta > 0\) , let \(\mathrm{N}_{\delta}\) be the set of ordered pairs \((x,y)\in W'\times W'\) such that \(\| x - y\| \leqslant \delta\) . The \(\mathrm{N}_{\delta}\) form a fundamental system of entourages in \(W'\) . We write \(W = \psi(W')\) . Let \(\mathrm{M}_{\delta}\) be the set of ordered pairs \((u,v)\in W\times W\) such that \(\| \phi (u v^{-1})\| \leqslant \delta\) . The \(\mathrm{M}_{\delta}\) form a fundamental system of entourages in \(W\) with the right uniform structure. But relation (2) proves that

\[
\mathrm{N} _ {\delta} \subset (\phi \times \phi) (\mathrm{M} _ {a ^ {- 1} \delta}), \quad (\phi \times \phi) (\mathrm{M} _ {\delta}) \subset \mathrm{N} _ {b \delta}
\]

and hence W has the property of the lemma.

PROPOSITION 1. A Lie group is a complete metrizable topological group.

Since \(e\) admits an open neighbourhood homeomorphic to an open ball of a normed space, \(e\) admits a countable fundamental system of neighbourhoods whose intersection is \(\{e\}\) . Hence G is metrizable (General Topology, Chapter III, § 1, Corollary to Proposition 2 and Chapter IX, § 3, Proposition 1). By Lemma 1, there exists a neighbourhood of \(e\) which is complete under the right uniform structure and hence G is complete (General Topology, Chapter III, § 3, Proposition 4.

PROPOSITION 2. Let G be a Lie group.

\begin{enumerate}
\def\labelenumi{(\roman{enumi})}
\item
  If \(\mathbf{K} = \mathbf{R}\) or \(\mathbf{C}\) , \(\mathbf{G}\) is locally connected.
\item
  If \(\mathbf{K}\) is distinct from \(\mathbf{R}\) and \(\mathbf{C}\) , \(\mathbf{G}\) is zero-dimensional (General Topology, Chapter IX, § 6, Definition 5).
\item
  Suppose that K is locally compact. For G to be locally compact, it is necessary and sufficient that G be finite-dimensional.
\item
  If \(\mathbf{G}\) is generated by a subspace whose topology admits a countable base, then the topology on \(\mathbf{G}\) admits a countable base.
\end{enumerate}

Let U be a neighbourhood of \(e\) . There exists an open neighbourhood \(\mathbf{U}_1\) of \(e\) contained in U and homeomorphic to an open ball of a normed space E over K. If \(\mathbf{K} = \mathbf{R}\) or \(\mathbf{C}\) , \(\mathbf{U}_1\) is connected, which proves (i). Suppose that K is ultrametric. There exists a neighbourhood \(\mathbf{U}_2\) of \(e\) which is closed in G and such that \(\mathbf{U}_2 \subset \mathbf{U}_1\) . Then there exists a neighbourhood \(\mathbf{U}_3\) of \(e\) such that \(\mathbf{U}_3 \subset \mathbf{U}_2\) and \(\mathbf{U}_3\) is open and closed relative to \(\mathbf{U}_1\) . Then \(\mathbf{U}_3\) is closed relative to \(\mathbf{U}_2\) and hence G and open relative to \(\mathbf{U}_1\) and hence G. This proves (ii). For G to be locally compact, it is necessary and sufficient that E be locally compact; if K is locally compact, this amounts to saying that E is finite-dimensional (Topological Vector Spaces, Chapter I, § 2, Theorem 3), whence (iii). Suppose that G is generated by a subset V and let \(W = V \cup V^{-1}\) ; then

\[
G = W \cup W ^ {2} \cup W ^ {3} \cup \dots ;
\]

if there exists a sequence dense in V, we see that there exists a sequence dense in G and, as G is metrizable (Proposition 1), the topology on G admits a countable base.

COROLLARY. If K = R or C and G is connected and finite-dimensional, then G is locally connected and locally compact and its topology admits a countable base.

Lemma 2. Let X be a manifold of class \(\mathbf{C}^r\) , \(e\) a point of X, U and V open neighbourhoods of \(e\) and \(m\) a mapping of class \(\mathbf{C}^r\) of \(\mathrm{U} \times \mathrm{U}\) into X satisfying the following conditions:

\begin{enumerate}
\def\labelenumi{(\alph{enumi})}
\item
  \(m(e,x) = m(x,e) = x\) for all \(x\in \mathbf{U}\) ;
\item
  \(\mathrm{V} \subset \mathrm{U}\) , \(m(\mathrm{V} \times \mathrm{V}) \subset \mathrm{U}\) and \(m(m(x, y), z) = m(x, m(y, z))\) for all \(x, y, z\) in \(\mathrm{V}\) .
\end{enumerate}

Then there exists an open neighbourhood W of \(e\) in V and an automorphism \(\theta\) of the manifold W such that \(\theta(e) = e\) , \(\theta(\theta(x)) = x\) and \(m(x, \theta(x)) = m(\theta(x), x) = e\) for all \(x \in W\) .

\(m(e,y) = y\) for all \(y \in \mathrm{U}\) and hence, by the implicit function theorem, there exists an open neighbourhood \(\mathrm{W}_1\) of \(e\) in \(\mathrm{V}\) and a mapping \(\theta_1\) of class \(C^r\) of \(\mathrm{W}_1\) into \(\mathrm{V}\) such that \(\theta_1(e) = e\) , \(m(x,\theta_1(x)) = e\) for all \(x \in \mathrm{W}_1\) . Similarly, there exists an open neighbourhood \(\mathrm{W}_2\) of \(e\) in \(\mathrm{V}\) and a mapping \(\theta_2\) of class \(C^r\) of \(\mathrm{W}_2\) into \(\mathrm{V}\) such that \(\theta_2(e) = e\) , \(m(\theta_2(x),x) = e\) for all \(x \in \mathrm{W}_2\) . For \(x \in \mathrm{W}_1 \cap \mathrm{W}_2\) .

\[
\begin{array}{r l} \theta_ {2} (x) & = m (\theta_ {2} (x), e) = m (\theta_ {2} (x), m (x, \theta_ {1} (x))) \\ & = m (m (\theta_ {2} (x), x), \theta_ {1} (x)) = m (e, \theta_ {1} (x)) = \theta_ {1} (x). \end{array}
\]

Let \(\theta(x)\) be the common value of \(\theta_1(x)\) and \(\theta_2(x)\) for \(x \in W_1 \cap W_2\) . Let \(W\) be the set of \(x \in W_1 \cap W_2\) such that \(\theta(x) \in W_1 \cap W_2\) . The set \(W\) is open. For \(x \in W\) ,

\[
\theta (\theta (x)) = m (m (x, \theta (x)), \theta (\theta (x))) = m (x, m (\theta (x), \theta (\theta (x)))) = m (x, e) = x
\]

and hence \(\theta(x) \in W\) . We see that \(\theta \mid W\) defines an automorphism of the manifold \(W\) .

PROPOSITION 3. Let \(\mathbf{X}\) be an analytic manifold and \(m\) an analytic associative law of composition on \(\mathbf{X}\) admitting an identity element. The set \(\mathbf{G}\) of invertible elements of \(\mathbf{X}\) is open in \(\mathbf{X}\) and \(\mathbf{G}\) is a Lie group with \(m \mid (\mathbf{G} \times \mathbf{G})\) and the manifold structure induced by that on \(\mathbf{X}\) .

By Lemma 2, G is a neighbourhood of the identity element. For all \(g \in \mathbb{G}\) , the mapping \(x \mapsto m(g, x)\) is an automorphism of the manifold X. Hence the image of G under this mapping is a neighbourhood of \(g\) , obviously contained in G. Therefore G is open in X. Clearly conditions \((\mathrm{GL}_1)\) and \((\mathrm{GL}_2)\) hold. Condition \((\mathrm{GL}_3)\) holds by Lemma 2.

